% Copyright (c) 2022-2026 MKM Research Labs. Licensed under the MIT License; see LICENSE.
%
% MKM-SG-001 Storm-Gauge Forward Model -- model documentation v2.0, the document of record
% filed on the model as Core Documentation. It supersedes v1.3 (11-Mar-2026). The Wolfram
% benchmark of v1.3 (MKM-SG-001_wolfram_benchmark v0.1) found the code computed its
% equations, but that a storm never reached the levels the calibration named, that the
% response did not behave as described, and that nothing in production used the model.
% v2.0 changes the model and describes it, and where it is used, as they are. Its tables
% are inlined, generated from the code.
%
% Build: it inputs the house header, which lives in PhysicalRisk. Copy
% PhysicalRisk/docs/models/shared beside this folder as ../shared, then run
% pdflatex three times (the third settles "Page n of N").
\documentclass[11pt,a4paper]{article}
\newcommand{\doctitle}{Storm-Gauge Forward Model}
\newcommand{\docsubtitle}{Spatial Decay and Intensity-to-Level Conversion}
\newcommand{\docversion}{2.0}
\newcommand{\docdate}{04-November-2026}
\newcommand{\docauthor}{David K Kelly}
\input{../shared/mkm_header}
% T1 so the underscores in code identifiers print as underscores.
\usepackage[T1]{fontenc}
\usepackage{lmodern}
% The shared history table sets every column to its natural width; wrap the description.
\renewcommand{\mkmhistorypage}[1]{%
    \section{Document History}
    \begin{table}[ht]
        \centering
        \begin{tabular}{l|l|p{9cm}|l}
            \textbf{Date} & \textbf{Version} & \textbf{Description} & \textbf{Author} \\
            \hline
            #1
        \end{tabular}
    \end{table}
}
\newcommand{\sgSteadyAlert}{38.7}
\newcommand{\sgSteadyWarning}{64.7}
\newcommand{\sgSteadySevere}{86.0}
\newcommand{\sgStormAlert}{40.9}
\newcommand{\sgStormWarning}{68.2}
\newcommand{\sgStormSevere}{90.7}
\newcommand{\sgLag}{2.5}
\newcommand{\sgGaugeMs}{0.3}
\newcommand{\sgAllMs}{8}
\newcommand{\sgGammaUp}{14.9}
\newcommand{\sgGammaDown}{12.9}
\newcommand{\sgGammaLo}{10}
\newcommand{\sgGammaHi}{20}
\begin{document}
\sloppy
\mkmtitlepage
\mkmlegalpage
\tableofcontents
\clearpage

% ===================================================================
% 1. EXECUTIVE SUMMARY
% ===================================================================
\mkmexecsummary

The Storm-Gauge Forward Model maps a storm's track, intensity profile and footprint to a
water level time series at each river gauge. It has three stages:
\begin{enumerate}[nosep]
    \item \textbf{Spatial proximity:} the storm's position, interpolated linearly along
          its track at each time step, and its Haversine distance from the gauge.
    \item \textbf{Intensity attenuation:} a Gaussian, exponential or linear decay
          kernel applied to the storm's intensity.
    \item \textbf{Intensity-to-level conversion:} a power law gives each intensity a
          target level, and the gauge approaches it after a lag, rising and falling
          with separate time constants.
\end{enumerate}

\paragraph{This is version 2.0, a revised model.} The Wolfram benchmark of v1.3
(MKM-SG-001\_wolfram\_benchmark v0.1, 04-Nov-2026) found that the code computed its
documented equations exactly. What the document said about them did not hold. v2.0
changes the model in three ways:
\begin{itemize}[nosep]
    \item \textbf{The rise has its own time constant,} $\tau_{\text{rise}}=1$\,h. In
          v1.3 the rise rate was tied to the fall, at $3\Delta t/\tau_{\text{decay}}$,
          a 4-hour rise. The gauge reached only three-quarters of its target before the
          storm passed. A storm therefore needed a peak of 56 to reach the flood alert,
          where the calibration said 40, and never reached severe. Doubling
          $\tau_{\text{decay}}$ lowered the peak. Now a storm crossing the gauge reaches
          alert at peak \sgStormAlert, warning at \sgStormWarning{} and severe at
          \sgStormSevere{} (\S\ref{sec:intensity_to_level}). $\tau_{\text{decay}}$ governs
          only the fall.
    \item \textbf{The run includes the recession.} v1.3 stopped at the storm's end,
          with the gauge still high, so durations and accumulation were cut short.
          v2.0 runs five fall constants further, with no forcing.
    \item \textbf{An unread input is removed.} \texttt{intensity\_to\_level\_scale} was
          stored and never read. The constants move to \texttt{config.models} (R1).
\end{itemize}

\paragraph{Where the model is used.} v1.3 described the model as feeding
\texttt{gaugets/}, the hazard curves and the flood risk model. It does not, and did not.
Nothing in the platform imports it. \texttt{gaugets/} is written by a scripted flood wave
(\texttt{port/rand/shared/gauge/gaugets\_random}). The model is a standalone forward
model, used for analysis and as the reference for the gauge response that MKM-SS-001
re-implements (\S\ref{sec:scope}).

The model is implemented in Python with only the standard library's \texttt{math}
module, so it is portable and easy to audit.

% ===================================================================
% 2. MODEL PURPOSE AND SCOPE
% ===================================================================
\mkmpurpose

\subsection{Purpose}
The model computes the water level response at each gauge in a catchment to a
parameterised storm. It supports:
\begin{itemize}[nosep]
    \item analysis of how a gauge responds to storm characteristics: intensity, track
          proximity, footprint width and kernel shape;
    \item the spatial pattern of flooding across a catchment's gauges for one storm;
    \item a reference for the lagged, asymmetric gauge response that other models
          re-implement.
\end{itemize}

\subsection{Scope}
\label{sec:scope}
\begin{itemize}[nosep]
    \item \textbf{In scope:} track interpolation; Haversine distance; the decay kernels;
          the intensity-to-level transfer; the lagged rise and fall; and summary
          statistics (peak level, durations above each stage, accumulation).
    \item \textbf{Out of scope:} storm generation (\texttt{storm\_factory.py} builds
          simple storms for the model's own use); hydrodynamic solvers; property
          depth--damage; PRS pricing.
    \item \textbf{Upstream:} \texttt{Storm} objects, from \texttt{create\_storm} or
          built directly; \texttt{GaugeConfig} objects, from the gauge portfolio CDM
          through \texttt{GaugeConfig.from\_gauge\_portfolio}.
    \item \textbf{Downstream: none in production.} No module in \texttt{src},
          \texttt{tools} or \texttt{scripts} imports the package. The gauge time series
          in \texttt{gaugets/} come from \texttt{port/src/gauge/gaugets.py}, which
          draws a scripted wave from \texttt{gaugets\_random}. MKM-SS-001's sequence
          response (\texttt{storm\_multi/models/sequence\_response.py}) re-implements
          a lagged, asymmetric rise and fall of the same kind; it does not call this
          code.
\end{itemize}

\subsection{Model Classification}
Under the firm's model risk management framework this is a \emph{Tier~2 quantitative
model}. It is deterministic given its inputs. Because no production model consumes its
output, an error in it cannot at present reach a price, a risk figure or a report. That
would change if it were wired into \texttt{gaugets/}, and this document would have to
change with it.

% ===================================================================
% 3. MATHEMATICAL FRAMEWORK
% ===================================================================
\mkmframework

\subsection{Haversine Distance}
\label{sec:haversine}

The distance between a gauge at $(\phi_g, \lambda_g)$ and a storm position
$(\phi_s, \lambda_s)$ is computed on a sphere of radius $R = 6{,}371$\,km:

\begin{equation}
    a = \sin^2\!\left(\frac{\Delta\phi}{2}\right)
      + \cos\phi_g \,\cos\phi_s \,\sin^2\!\left(\frac{\Delta\lambda}{2}\right)
    \label{eq:haversine_a}
\end{equation}

\begin{equation}
    d = 2R \,\arctan\!\left(\frac{\sqrt{a}}{\sqrt{1-a}}\right)
    \label{eq:haversine_d}
\end{equation}

\noindent where $\Delta\phi = \phi_s - \phi_g$ and $\Delta\lambda = \lambda_s - \lambda_g$,
in radians. The code evaluates Eq.~\ref{eq:haversine_d} with \texttt{math.atan2}, which
equals $\arcsin\sqrt{a}$ without its loss of accuracy as $a\to1$.

\subsection{Storm Track Interpolation}
\label{sec:track_interp}

A storm is an ordered sequence of track points $\{(\lambda_i, \phi_i, t_i, I_i)\}_{i=1}^{N}$:
longitude, latitude, hours since the storm started, and intensity on a $[0, 100]$ scale.
At time $t$ the model finds the bracketing points with
$t_{\text{prev}} \leq t < t_{\text{next}}$ and the fraction

\begin{equation}
    \alpha = \operatorname{clamp}\!\left(
        \frac{t - t_{\text{prev}}}{t_{\text{next}} - t_{\text{prev}}},\; 0,\; 1
    \right)
    \label{eq:interp_frac}
\end{equation}

\noindent and interpolates linearly:

\begin{align}
    \lambda_s(t) &= \lambda_{\text{prev}} + \alpha\,(\lambda_{\text{next}} - \lambda_{\text{prev}}) \label{eq:interp_lon} \\
    \phi_s(t) &= \phi_{\text{prev}} + \alpha\,(\phi_{\text{next}} - \phi_{\text{prev}}) \label{eq:interp_lat} \\
    I_s(t) &= I_{\text{prev}} + \alpha\,(I_{\text{next}} - I_{\text{prev}}) \label{eq:interp_intensity}
\end{align}

\noindent Before the first track point the storm is at that point; after the last, at the
last. After the storm's duration $T$ its intensity is zero (\S\ref{sec:temporal_dynamics}).

\subsection{Spatial Decay Kernels}
\label{sec:decay_kernels}

The storm's footprint radius $F$ (km) normalises the distance,

\begin{equation}
    \tilde{d} = \frac{d}{F},
    \label{eq:d_norm}
\end{equation}

\noindent and one of three kernels attenuates the intensity:

\subsubsection{Gaussian Kernel}
\begin{equation}
    k_{\text{Gauss}}(\tilde{d}\,;\,\sigma) =
        \exp\!\left(-\frac{\tilde{d}^{\,2}}{2\sigma^2}\right)
    \label{eq:gaussian_kernel}
\end{equation}

\noindent with bandwidth $\sigma$ (default 0.5). It is the default kernel. At
$\tilde{d} = \sigma$ the intensity is 60.7\% of its on-track value, and at $2\sigma$
it is 13.5\%.

\subsubsection{Exponential Kernel}
\begin{equation}
    k_{\text{Exp}}(\tilde{d}\,;\,\lambda) =
        \exp\!\left(-\frac{\tilde{d}}{\lambda}\right)
    \label{eq:exponential_kernel}
\end{equation}

\noindent with decay length $\lambda$ (default 0.5). It is sharper at the centre than
the Gaussian, and has heavier tails.

\subsubsection{Linear Kernel}
\begin{equation}
    k_{\text{Lin}}(\tilde{d}\,;\,r) =
        \max\!\left(0,\; 1 - \frac{\tilde{d}}{r}\right)
    \label{eq:linear_kernel}
\end{equation}

\noindent with cut-off $r$ (default 0.5). Intensity is zero at and beyond $\tilde{d}=r$.

\paragraph{Boundary condition.} For every kernel $k(0) = 1$: a gauge on the track
receives the full intensity. The Gaussian and exponential kernels tend to zero with
distance; the linear kernel reaches it.

\subsection{Intensity at Gauge}
\label{sec:intensity_at_gauge}

The intensity a gauge experiences at time $t$ is

\begin{equation}
    I_{\text{gauge}}(t) = I_s(t) \times k\!\left(\tilde{d}(t)\right),
    \label{eq:intensity_at_gauge}
\end{equation}

\noindent with $I_s(t)$ from Eq.~\ref{eq:interp_intensity} and $\tilde{d}(t)$ the
normalised distance to the storm's position at $t$. Intensity and footprint are assumed
separable (\S8, item~3).

\subsection{Intensity-to-Level Conversion}
\label{sec:intensity_to_level}

Each intensity has a target level. For a gauge with base level $L_{\text{base}}$,
historical high $L_{\text{high}}$ (or $1.1 \times L_{\text{severe}}$ if it has none) and
sensitivity $s$:

\begin{equation}
    L_{\text{range}} = L_{\text{high}} - L_{\text{base}}
    \label{eq:level_range}
\end{equation}

\begin{equation}
    \Delta L(I) = \left(\frac{I}{100}\right)^{\!\gamma} \times L_{\text{range}} \times s,
    \qquad \Delta L(I) = 0 \text{ for } I \le 0,
    \label{eq:power_law}
\end{equation}

\noindent with $\gamma = 0.8$. The sub-linear exponent gives moderate intensities a
proportionally larger target than a linear map would.

\paragraph{The calibration bands.} At Table~\ref{tab:gauge_params}'s default stages the
target reaches each stage at a fixed intensity, whatever the alert level: flood alert at
$I=\sgSteadyAlert$, flood warning at \sgSteadyWarning, severe at \sgSteadySevere{} and the
historical high at 100. A target is only reached if the intensity is held for long
enough. A storm passes, so it must peak higher. A default storm
(Table~\ref{tab:storm_params}) crossing directly over the gauge reaches:
\begin{itemize}[nosep]
    \item flood alert at peak intensity \sgStormAlert;
    \item flood warning at \sgStormWarning;
    \item severe at \sgStormSevere.
\end{itemize}
These storm figures are the ones v2.0's rise constant was chosen to give. The design
intent is alert near 40, warning in 60--70 and severe at 85 or above
(Table~\ref{tab:storm_response}). A gauge with its own historical high has its own bands.

\subsection{Temporal Response Dynamics}
\label{sec:temporal_dynamics}

The water level at step $i$ is

\begin{equation}
    L(t_i) = L_{\text{base}} + C(t_i),
    \label{eq:gauge_level}
\end{equation}

\noindent where $C$ is the contribution. With time step $\Delta t$ (0.5\,h) and lag
$\tau_{\text{lag}}$ (2\,h), the target at step $i$ is $\Delta L$ of the lagged intensity

\begin{equation}
    I_{\text{lag}}(t_i) = I_{\text{gauge}}\!\left(t_{\max(0,\; i - \lfloor \tau_{\text{lag}}/\Delta t \rfloor)}\right),
    \label{eq:lagged_intensity}
\end{equation}

\noindent and the contribution approaches it with a rise constant $\tau_{\text{rise}}$
(1\,h) and a fall constant $\tau_{\text{decay}}$ (12\,h):

\begin{equation}
    C(t_i) = C(t_{i-1}) + \bigl[\Delta L_{\text{target}} - C(t_{i-1})\bigr] \times
    \begin{cases}
        \min\!\left(1,\; \Delta t/\tau_{\text{rise}}\right) & \text{if } \Delta L_{\text{target}} > C(t_{i-1}) \\[0.5em]
        \min\!\left(1,\; \Delta t/\tau_{\text{decay}}\right) & \text{otherwise,}
    \end{cases}
    \label{eq:contribution_update}
\end{equation}

\noindent from $C=0$. The gauge rises twelve times faster than it falls.

\paragraph{The run.} The grid is $t_i = i\,\Delta t$ for
$i = 0, \ldots, \lfloor (T + 5\tau_{\text{decay}})/\Delta t \rfloor$. After the storm's
duration $T$ there is no forcing. The five fall constants that follow let the
contribution fall to $e^{-5}$, under 1\%, of where it started falling.

\paragraph{Lag between the peaks.} Peak level follows peak intensity by
$\tau_{\text{lag}}$ plus the time the rising level takes to meet the falling target.
For a default storm that is \sgLag\,h (Table~\ref{tab:storm_response}).

\subsection{Summary Statistics}
\label{sec:summary_stats}

Over the whole run, recession included:

\begin{itemize}[nosep]
    \item \textbf{Peak level} $L_{\text{peak}} = \max_i L(t_i)$, and its time.
    \item \textbf{Peak exceedance} $E_{\text{peak}} = L_{\text{peak}} - L_{\text{alert}}$,
          negative if the gauge does not reach alert.
    \item \textbf{Duration above each stage} $\theta$ (alert, warning, severe):
          \begin{equation}
              D_{\theta} = \sum_{i:\, L(t_i) \geq \theta} \Delta t
              \label{eq:duration}
          \end{equation}
    \item \textbf{Accumulation} above the alert level, in metre-hours:
          \begin{equation}
              A = \sum_{i} \max\!\bigl(0,\; L(t_i) - L_{\text{alert}}\bigr)\,\Delta t
              \label{eq:accumulation}
          \end{equation}
    \item \textbf{Peak intensity} $I_{\text{peak}} = \max_i I_{\text{gauge}}(t_i)$,
          and its time.
\end{itemize}

% ===================================================================
% 4. INPUT PARAMETERS AND CALIBRATION
% ===================================================================
\mkminputs

\subsection{Storm Parameters}

\begin{table}[H]
    \centering
    \caption{Storm input parameters (defaults of \texttt{create\_storm})}
    \label{tab:storm_params}
    \begin{tabular}{llll}
        \toprule
        \textbf{Parameter} & \textbf{Symbol} & \textbf{Units} & \textbf{Default} \\
        \midrule
        Track points & $\{(\lambda_i, \phi_i, t_i, I_i)\}$ & deg, h, --- & 25 points \\
        Peak intensity & $I_{\text{peak}}$ & 0--100 & 50.0 \\
        Footprint radius & $F$ & km & 40.0 \\
        Duration & $T$ & hours & 24.0 \\
        Decay kernel & --- & enum & Gaussian \\
        Decay parameter & $\sigma$, $\lambda$ or $r$ & --- & 0.5 \\
        Intensity profile & --- & enum & Triangular \\
        \bottomrule
    \end{tabular}
\end{table}

\noindent The peak intensity and profile shape the track's intensities when
\texttt{create\_storm} builds it; the model reads only the track. A \texttt{Storm} also
carries a speed, which the model does not read.

\subsection{Gauge Configuration Parameters}

\begin{table}[H]
    \centering
    \caption{Gauge configuration parameters (defaults of \texttt{from\_gauge\_portfolio})}
    \label{tab:gauge_params}
    \begin{tabular}{llll}
        \toprule
        \textbf{Parameter} & \textbf{Symbol} & \textbf{Units} & \textbf{Default} \\
        \midrule
        Base level & $L_{\text{base}}$ & m & $0.35 \times L_{\text{alert}}$ \\
        Flood alert & $L_{\text{alert}}$ & m & 4.0 \\
        Flood warning & $L_{\text{warn}}$ & m & $1.33 \times L_{\text{alert}}$ \\
        Severe warning & $L_{\text{severe}}$ & m & $1.58 \times L_{\text{alert}}$ \\
        Historical high & $L_{\text{high}}$ & m & optional \\
        Sensitivity & $s$ & --- & 1.0 \\
        \bottomrule
    \end{tabular}
\end{table}

\subsection{Model Hyperparameters}

\begin{table}[H]
    \centering
    \caption{Model hyperparameters (\texttt{config.models}, \texttt{\_stormgauge.py})}
    \label{tab:model_params}
    \begin{tabular}{llll}
        \toprule
        \textbf{Parameter} & \textbf{Symbol} & \textbf{Units} & \textbf{Default} \\
        \midrule
        Time resolution & $\Delta t$ & hours & 0.5 \\
        Response lag & $\tau_{\text{lag}}$ & hours & 2.0 \\
        Rise time constant & $\tau_{\text{rise}}$ & hours & 1.0 \\
        Fall time constant & $\tau_{\text{decay}}$ & hours & 12.0 \\
        Recession run-on & --- & $\tau_{\text{decay}}$ & 5 \\
        Power-law exponent & $\gamma$ & --- & 0.8 \\
        Historical high, where none & --- & $\times L_{\text{severe}}$ & 1.1 \\
        Earth radius & $R$ & km & 6{,}371 \\
        \bottomrule
    \end{tabular}
\end{table}

\subsection{Calibration Approach}

The parameters are design choices, not fitted values.
\begin{itemize}[nosep]
    \item $\gamma = 0.8$ and the stage multiples of Table~\ref{tab:gauge_params} set
          the steady-state bands of \S\ref{sec:intensity_to_level}.
    \item $\tau_{\text{rise}} = 1$\,h was chosen in v2.0 so that a storm crossing the
          gauge reaches the bands those choices intend.
    \item $\tau_{\text{lag}}=2$\,h and $\tau_{\text{decay}}=12$\,h are unchanged from
          v1.3.
\end{itemize}
v1.3 described the parameters as calibrated against Environment Agency records for the
Thames. No such calibration is held in the repository, and v2.0 does not repeat the
claim. Calibrating them to observed hydrographs is the first item of \S8.

The gauge sensitivity $s$ (default 1.0) can be set per gauge, for catchment
characteristics such as urbanisation and soil.

\subsection{Data Formats}

\subsubsection{Storm Track Data}
A storm's track is a list of \texttt{TrackPoint} objects:
\begin{lstlisting}[language=Python]
@dataclass
class TrackPoint:
    longitude: float    # Degrees (WGS84)
    latitude: float     # Degrees (WGS84)
    time_hours: float   # Hours since storm start
    intensity: float    # 0-100 intensity scale
\end{lstlisting}

\subsubsection{Gauge Configuration}
\texttt{GaugeConfig.from\_gauge\_portfolio()} reads a gauge's stages and coordinates from
its gauge portfolio CDM entry, with the defaults of Table~\ref{tab:gauge_params} where a
stage is missing.

% ===================================================================
% 5. IMPLEMENTATION DETAILS
% ===================================================================
\mkmimplementation

\subsection{Software Architecture}

The package \texttt{src/models/stormgauge/} holds:
\begin{itemize}[nosep]
    \item \texttt{data\_structures.py}: the enumerations \texttt{DecayKernel} and
          \texttt{IntensityProfile}, and the dataclasses \texttt{TrackPoint},
          \texttt{Storm}, \texttt{GaugeConfig} and \texttt{GaugeResponse};
    \item \texttt{forward\_model/}: \texttt{StormGaugeModel} (\texttt{\_core.py}), with
          distance, decay, interpolation and the power law in \texttt{\_intensity.py}
          and the response in \texttt{\_response.py};
    \item \texttt{storm\_factory.py}: \texttt{create\_storm} and
          \texttt{load\_gauges\_from\_portfolio}.
\end{itemize}
Its parameters are in \texttt{config/models/\_stormgauge.py}.

\subsection{Computational Flow}

\texttt{StormGaugeModel.compute\_response(storm, gauge)}:
\begin{enumerate}[nosep]
    \item builds the grid of \S\ref{sec:temporal_dynamics}, from 0 to
          $T + 5\tau_{\text{decay}}$;
    \item computes the intensity at the gauge at each step (Eqs.~3--11), zero after $T$;
    \item applies the lagged target and Eq.~\ref{eq:contribution_update};
    \item extracts the statistics of \S\ref{sec:summary_stats};
    \item returns a \texttt{GaugeResponse} with the series and the statistics.
\end{enumerate}
\texttt{compute\_all\_responses(storm, gauges)} runs each gauge in turn. The cost is
$\mathcal{O}(N_t \cdot N_{\text{track}})$ per gauge.

\subsection{Numerical Considerations}

\begin{itemize}[nosep]
    \item The Haversine uses \texttt{math.atan2} (\S\ref{sec:haversine}).
    \item $\alpha$ is clamped to $[0, 1]$, so the storm never leaves its track.
    \item Both smoothing rates are clamped by $\min(1,\cdot)$, so a large step cannot
          overshoot its target.
    \item Arithmetic is 64-bit floating point. A kernel value below the smallest
          double, far outside the footprint, is zero.
\end{itemize}

\subsection{Dependencies}

The Python standard library (\texttt{math}, \texttt{dataclasses}, \texttt{enum},
\texttt{typing}) and \texttt{config.models}.

% ===================================================================
% 6. SENSITIVITY ANALYSIS
% ===================================================================
\mkmsensitivity

\begin{table}[H]
    \centering
    \caption{Spatial decay factor by kernel type (footprint 100\,km, decay parameter 1.0).}
    \label{tab:sens_decay_kernel}
    \begin{tabular}{rrrr}
        \toprule
        \textbf{Distance (km)} & \textbf{Gaussian} & \textbf{Exponential} & \textbf{Linear} \\
        \midrule
        0 & 1.00 & 1.00 & 1.00 \\
        25 & 0.9692 & 0.7788 & 0.7500 \\
        50 & 0.8825 & 0.6065 & 0.5000 \\
        75 & 0.7548 & 0.4724 & 0.2500 \\
        100 & 0.6065 & 0.3679 & 0.0000 \\
        150 & 0.3247 & 0.2231 & 0.0000 \\
        200 & 0.1353 & 0.1353 & 0.0000 \\
        \bottomrule
    \end{tabular}
\end{table}

\begin{table}[H]
    \centering
    \caption{Target contribution $\Delta L$ (m, Eq.~\ref{eq:power_law}) by intensity and gauge sensitivity ($L_{\text{base}}=1.5$, $L_{\text{high}}=6.0$\,m).}
    \label{tab:sens_intensity_level}
    \begin{tabular}{rrrrrr}
        \toprule
        \textbf{$s$} & \textbf{$I=20$} & \textbf{$I=40$} & \textbf{$I=60$} & \textbf{$I=80$} & \textbf{$I=100$} \\
        \midrule
        0.8 & 0.9930 & 1.73 & 2.39 & 3.01 & 3.60 \\
        1.0 & 1.24 & 2.16 & 2.99 & 3.76 & 4.50 \\
        1.2 & 1.49 & 2.59 & 3.59 & 4.52 & 5.40 \\
        \bottomrule
    \end{tabular}
\end{table}

\subsection{Decay Kernel Selection}

The kernel matters for storms that do not pass over the gauge. One footprint radius
away ($\tilde{d} = 1.0$):

\begin{table}[H]
    \centering
    \caption{Decay factor at $\tilde{d} = 1.0$ for each kernel (default parameters)}
    \label{tab:decay_comparison}
    \begin{tabular}{lcc}
        \toprule
        \textbf{Kernel} & \textbf{$k(1.0)$} & \textbf{Relative intensity} \\
        \midrule
        Gaussian ($\sigma = 0.5$) & $\exp(-2) \approx 0.135$ & 13.5\% \\
        Exponential ($\lambda = 0.5$) & $\exp(-2) \approx 0.135$ & 13.5\% \\
        Linear ($r = 0.5$) & $\max(0, 1-2) = 0.0$ & 0.0\% \\
        \bottomrule
    \end{tabular}
\end{table}

\noindent The Gaussian and exponential kernels coincide here for $\sigma = \lambda = 0.5$
and nowhere else. The linear kernel gives nothing beyond $\tilde{d} = r = 0.5$.

\subsection{Power-Law Exponent}

\begin{table}[H]
    \centering
    \caption{Normalised level contribution $\Delta L / L_{\text{range}}$ at $I = 50$ for varying $\gamma$}
    \label{tab:gamma_sensitivity}
    \begin{tabular}{lc}
        \toprule
        \textbf{$\gamma$} & \textbf{$\Delta L / (s \cdot L_{\text{range}})$} \\
        \midrule
        0.5 & $(0.5)^{0.5} \approx 0.707$ \\
        0.8 (default) & $(0.5)^{0.8} \approx 0.574$ \\
        1.0 & $(0.5)^{1.0} = 0.500$ \\
        1.5 & $(0.5)^{1.5} \approx 0.354$ \\
        \bottomrule
    \end{tabular}
\end{table}

\noindent A 25\% change in $\gamma$ moves the target at $I=50$ by $+\sgGammaUp$\% (down
to 0.6) and $-\sgGammaDown$\% (up to 1.0). Across $I=40$--60 the change is
\sgGammaLo--\sgGammaHi\%. The effect grows as the intensity falls.

\subsection{Storm Response}

\begin{table}[H]
    \centering
    \caption{Response of a default storm (Table~\ref{tab:storm_params}) crossing the Thames at 51.5$^\circ$N, at a gauge with Table~\ref{tab:gauge_params}'s default stages ($L_{\text{alert}}=4$\,m) on the track and 0.1 and 0.2$^\circ$ north of it. Lag is peak level less peak intensity, in hours; durations in hours; accumulation in metre-hours.}
    \label{tab:storm_response}
    \begin{tabular}{rrrrrrrrr}
        \toprule
        \textbf{Peak} & \textbf{North ($^\circ$)} & \textbf{$I_{\text{peak}}$} & \textbf{$L_{\text{peak}}$ (m)} & \textbf{Lag} & \textbf{$D_{\text{alert}}$} & \textbf{$D_{\text{warn}}$} & \textbf{$D_{\text{severe}}$} & \textbf{$A$} \\
        \midrule
        40 & 0.0 & 40.0 & 3.96 & 2.5 & 0.0 & 0.0 & 0.0 & 0.00 \\
        40 & 0.1 & 34.3 & 3.66 & 2.5 & 0.0 & 0.0 & 0.0 & 0.00 \\
        40 & 0.2 & 21.6 & 2.96 & 2.5 & 0.0 & 0.0 & 0.0 & 0.00 \\
        60 & 0.0 & 60.0 & 4.94 & 2.5 & 10.5 & 0.0 & 0.0 & 6.20 \\
        60 & 0.1 & 51.4 & 4.52 & 2.5 & 8.0 & 0.0 & 0.0 & 2.68 \\
        60 & 0.2 & 32.3 & 3.56 & 2.5 & 0.0 & 0.0 & 0.0 & 0.00 \\
        80 & 0.0 & 80.0 & 5.85 & 2.5 & 14.5 & 6.5 & 0.0 & 16.17 \\
        80 & 0.1 & 68.5 & 5.33 & 2.5 & 12.5 & 1.0 & 0.0 & 10.24 \\
        80 & 0.2 & 43.1 & 4.11 & 2.5 & 3.5 & 0.0 & 0.0 & 0.29 \\
        100 & 0.0 & 100.0 & 6.72 & 2.5 & 17.5 & 10.5 & 5.0 & 27.36 \\
        100 & 0.1 & 85.7 & 6.10 & 2.5 & 15.5 & 8.0 & 0.0 & 19.27 \\
        100 & 0.2 & 53.9 & 4.65 & 2.5 & 8.5 & 0.0 & 0.0 & 3.63 \\
        \bottomrule
    \end{tabular}
\end{table}

\noindent The lag between the peaks is the same in every case, \sgLag\,h: the storms
share one profile, and the gauge's distance scales the forcing without reshaping it.
Off the track the storm is weaker: 0.2$^\circ$ (22\,km) north, a peak-100 storm
reaches alert but not warning.

\subsection{Time Constants}

\begin{table}[H]
    \centering
    \caption{Sensitivity to the time constants, on the peak-100 storm crossing the gauge; the others at their defaults.}
    \label{tab:time_constants}
    \begin{tabular}{lrrrrr}
        \toprule
        \textbf{Constant} & \textbf{Hours} & \textbf{$L_{\text{peak}}$ (m)} & \textbf{Lag (h)} & \textbf{$D_{\text{alert}}$ (h)} & \textbf{$A$ (m\,h)} \\
        \midrule
        $\tau_{\text{lag}}$ & 1 & 6.72 & 1.5 & 17.5 & 27.36 \\
        $\tau_{\text{lag}}$ & 2 & 6.72 & 2.5 & 17.5 & 27.36 \\
        $\tau_{\text{lag}}$ & 4 & 6.72 & 4.5 & 17.5 & 27.36 \\
        $\tau_{\text{rise}}$ & 0.5 & 6.95 & 2.0 & 18.5 & 30.66 \\
        $\tau_{\text{rise}}$ & 1 & 6.72 & 2.5 & 17.5 & 27.36 \\
        $\tau_{\text{rise}}$ & 2 & 6.31 & 3.0 & 16.0 & 21.53 \\
        $\tau_{\text{rise}}$ & 4 & 5.53 & 4.0 & 13.0 & 12.17 \\
        $\tau_{\text{decay}}$ & 6 & 6.72 & 2.5 & 13.5 & 21.38 \\
        $\tau_{\text{decay}}$ & 12 & 6.72 & 2.5 & 17.5 & 27.36 \\
        $\tau_{\text{decay}}$ & 24 & 6.72 & 2.5 & 26.0 & 38.07 \\
        \bottomrule
    \end{tabular}
\end{table}

\begin{itemize}[nosep]
    \item $\tau_{\text{lag}}$ moves the timing one for one and leaves the peak alone.
    \item $\tau_{\text{decay}}$ leaves the peak alone and lengthens the recession:
          doubling it from 12 to 24\,h takes the time above alert from 17.5 to
          26\,h.
    \item $\tau_{\text{rise}}$ sets how far the gauge gets before the storm passes.
          At 4\,h, v1.3's effective value, the peak-100 storm's peak is 1.2\,m lower.
\end{itemize}

\subsection{Footprint Radius}

The footprint $F$ scales the reach of the storm: doubling it doubles the distance at
which a given decay is reached. It decides which gauges a storm affects, and so drives
the spatial pattern of flooding across a catchment.

% ===================================================================
% 7. VALIDATION AND BACKTESTING
% ===================================================================
\mkmvalidation

\subsection{Unit Testing}

The package's tests (\texttt{tests/models/stormgauge/}) verify:
\begin{itemize}[nosep]
    \item the Haversine against known distances, and its symmetry;
    \item the kernels' boundary conditions: $k(0) = 1$, and the linear kernel zero for
          $\tilde{d} \geq r$;
    \item that a higher intensity gives a higher target, and zero intensity none;
    \item no forcing after the storm, and a completed recession;
    \item that $\tau_{\text{decay}}$ does not slow the rise, and lengthens the time above
          alert;
    \item that a held intensity brings the level to its Eq.~\ref{eq:power_law} target.
\end{itemize}

\subsection{Benchmark}

The Wolfram oracle benchmark (MKM-SG-001\_wolfram\_benchmark) checks the code against a
50-digit implementation of this document's equations and against Wolfram's own. It also
tests this document's figures and statements. Its v0.1, against v1.3, found the
shortcomings this version corrects.

\subsection{Backtesting}

No backtest is on record. v1.3 reported peak errors within $\pm0.3$\,m and timing
within $\pm1$\,h against Environment Agency records. It also reported cross-checks
against GEV fits. Nothing in the repository runs the model against gauge records, so
those figures cannot be reproduced, and v2.0 does not repeat them. A backtest against
observed Thames hydrographs is the condition of Q9.

% ===================================================================
% 8. MODEL LIMITATIONS AND KNOWN WEAKNESSES
% ===================================================================
\mkmlimitations

\begin{enumerate}[nosep]
    \item \textbf{Uncalibrated.} $\gamma$, $\tau_{\text{lag}}$, $\tau_{\text{rise}}$ and
          $\tau_{\text{decay}}$ are design choices. None has been fitted to observed
          hydrographs.
    \item \textbf{No antecedent conditions.} Soil moisture, groundwater and preceding
          rain are ignored: a storm always produces the same response. MKM-SS-001
          composites sequences of storms with partial recovery between them.
    \item \textbf{Separable intensity and decay.} The footprint's shape does not vary
          with intensity.
    \item \textbf{Static sensitivity.} $s$ does not vary with level or season.
    \item \textbf{Spherical Earth.} Against the WGS84 ellipsoid the sphere understates
          distances in the UK by up to 0.32\%, east--west, and by 0.06\% north--south.
    \item \textbf{One rise and one fall constant} for every gauge, whatever its
          catchment.
    \item \textbf{No tidal interaction} at estuarine or coastal gauges.
    \item \textbf{One footprint per storm,} where real storms have several bands.
    \item \textbf{Linear track interpolation} in latitude and longitude rather than
          along great circles.
\end{enumerate}

% ===================================================================
% 9. GOVERNANCE VALIDATION QUESTIONS
% ===================================================================
\mkmgovernance

\mkmvqitem{1}{Ronseal Test}{Does the model do what it says on the tin?}
Yes, as v2.0 describes it. It maps a storm to a water level series and summary statistics
at each gauge, and the Wolfram benchmark confirms it computes this document's equations.
v1.3 also said it supplied the platform's gauge time series, which it does not
(\S\ref{sec:scope}).

\mkmvqitem{2}{Accuracy and Stability}{Does the model provide timely, accurate and stable results?}
It is deterministic and numerically exact to double precision. Its accuracy against real
gauges is unknown: no backtest is on record (\S7.3).

\mkmvqitem{3}{Resource Usage}{Does the model consume excessive computational resources?}
No. A 48-hour storm with 25 track points takes \sgGaugeMs\,ms per gauge, recession
included, and 25 gauges take \sgAllMs\,ms.

\mkmvqitem{4}{Override Controls}{Are there controls over manual overrides?}
The per-gauge sensitivity $s$ comes from gauge configuration data. The other parameters
are in \texttt{config.models} under version control, and can be passed to
\texttt{StormGaugeModel} for analysis.

\mkmvqitem{5}{Output Consumption}{How is the output consumed for risk decision-making?}
It is not. No production model, report or screen consumes the model's output. It is used
for analysis, and as the reference shape for MKM-SS-001's sequence response.

\mkmvqitem{6}{Model Chaining}{Is the model's output used as input to another model?}
No. Its inputs are storms built by \texttt{create\_storm} or directly, and gauges from
the gauge portfolio CDM. MKM-SS-001 re-implements a response of the same kind; it does
not call this code.

\mkmvqitem{7}{Weaknesses}{Does the model have weaknesses and limitations of use?}
Yes, nine (\S8). The most material is that its parameters are uncalibrated.

\mkmvqitem{8}{Monitoring}{What type of monitoring is required?}
While nothing consumes its output, a re-run of the benchmark on any change to the package
or its parameters. If it is wired into production, the backtest of Q9 and an annual
comparison against Environment Agency records.

\mkmvqitem{9}{Overall Risk}{Is the overall model risk acceptable?}
Acceptable for its present use, which reaches no decision. Before its output is used in
production it needs calibration and a backtest against observed Thames hydrographs, and
this document must be revised to say where the output goes.

\clearpage

% ===================================================================
% 10. CHANGE HISTORY
% ===================================================================
\mkmchangelog

\textbf{v2.0 (04-Nov-2026).} After the Wolfram benchmark of v1.3:
\begin{itemize}[nosep]
    \item a rise time constant of its own, $\tau_{\text{rise}}=1$\,h, in place of
          $3\Delta t/\tau_{\text{decay}}$;
    \item the run continues $5\tau_{\text{decay}}$ past the storm's end, with no forcing,
          so durations and accumulation include the recession;
    \item \texttt{intensity\_to\_level\_scale}, never read, removed; the constants moved
          to \texttt{config.models};
    \item the scope, calibration, backtesting and governance answers corrected to what
          the model is and where it is used;
    \item the $\gamma$ sensitivity and the sphere's error restated as computed;
    \item Table~5 now shows sensitivity, and Tables~8--9 the storm response and the time
          constants.
\end{itemize}

% ===================================================================
% DOCUMENT HISTORY
% ===================================================================
\mkmhistorypage{13-Feb-2026 & 1.0 & Initial documentation & D. Kelly \\
18-Feb-2026 & 1.1 & Added Section 9: Governance Validation Questions (Chapter 5) & D. Kelly \\
26-Feb-2026 & 1.2 & Added Delta Engine model to inventory; updated models in scope & D. Kelly \\
11-Mar-2026 & 1.3 & Antecedent conditions partly addressed by MKM-SS-001 & D. Kelly \\
04-Nov-2026 & 2.0 & Revised model after the Wolfram benchmark: own rise constant, recession, unread scale removed; scope and governance corrected & D. Kelly \\}

\end{document}
